% ============================================================================
% October 16, 2026 pass: base-risk numbers moved to the constant-decay 2023
% state (decay_ship_20261014: table.md, qt_slide_numbers_decay.csv,
% slide33_main_values.json); quantum-tempo, credit-tenure and credit-for-
% parents framing as in the decay deck. Lower-risk rows and the earlier
% figures are still the two-step path (rescaled lower-risk decay reruns are
% on the Experiments side). Supply text untouched.
% October 13, 2026 build (agent draft). Policy section at the 7.9167
% production point: from-2023 GE transitions, each policy against its own
% no-policy path (base risk vs BASE, lower risk vs LR_BASE). Sources:
%   numbers: output/model/experiments/transitions_2023/7p9167_v6/table.md
%     (read Oct 13); housing behaviour: 7p9167_v6/housing_behavior/table.md
%     (= /mnt/project-files/housing_behavior_table_7p9167.md);
%   timing vs level: 7p9167_v6/notes/tempo_quantum_supm_vs_phi100.md;
%   long run and lower-risk construction: notes/results_evaluation_20261012/
%     evaluation.md sections 2A, 2F, 2I (checked against saved runs by the
%     Experiments thread, Oct 12);
%   starter home: 7p9167_v6/starter_home/starter_home_trap_by_age_7p9167_base.pdf
%     (values as read for the Oct 12 deck: ownership 22-25 46 -> 68%, 2-room
%     owners 8 -> 25%, births per 100 households 22.2 both);
%   PSID: notes/starter_home_trap_empirics_20261009.md, finding 3.
% Figures: tag-free copies (notag_7p9) used in the Oct 12 deck.
% The rebated property tax, transfers, fixed stock and eta = 1.5 arms are in
% Appendix app:policy-more (author Oct 12: no insurance and other policies in
% the main text for now).
% ============================================================================

\section{Housing Policy from 2023}
\label{sec:policy}


This section asks which housing policies can raise fertility in the 2023
economy. I consider two policies that are often proposed and differ in which
cost of family space they lower: an expansion of the housing supply, which
lowers the price of space and with it the per-period cost of a room, and
easier mortgage credit, which lowers the equity needed to buy. I then consider
two variants that sharpen the comparison: a supply expansion sized to raise
births on impact as much as credit does, and credit restricted to parents.

\subsection{Design of the experiments}
\label{sec:policy-design}

Every experiment starts from the 2023 state of the estimated transition of
Section~\ref{sec:quant-transition}: the same households, with the same
wealth, homes and children. In 2023 a policy is announced unexpectedly and
believed to be permanent. Households then have perfect foresight about the
paths of the house price, rents, the pension and the rebate, which clear their
markets at every date until the economy reaches a new stationary equilibrium.
I compare each policy with the path the economy follows from the same state
without the policy, so that the differences isolate the policy from the
ongoing transition. The preference for children stays at its 2023 value
throughout. I report three measures. \emph{Births on impact} are births in
2023--2026 relative to the no-policy path. \emph{Completed fertility by
cohort} is the number of children ever born to the cohort aged 18 to 21 in a
given year, relative to the same cohort without the policy. The \emph{long
view} follows the price, period fertility and the number of households to
2100.

The policies are as follows. \emph{Housing supply}: the supply curve
\eqref{eq:m-supply} shifts out by 10 percent from 2023 on, $H_0\to1.1H_0$,
with the same elasticity. \emph{Mortgage credit}: the financed share rises
from $\phi=0.80$ to 1, so buyers need no down payment, and owners who stay can
borrow up to the value of their home. \emph{Credit for parents}: the financed
share rises to 1 only for households with a child at home. \emph{Matched
supply}: the supply curve shifts out by 6.3 percent, a size chosen so that
births in 2023--2026 rise as much as under full mortgage credit.

I evaluate each policy twice: at the estimated earnings risk and at lower
earnings risk, with the standard deviation of the four-year earnings shock
lowered from 0.48 to 0.35 from 2023 on. Each policy under lower risk is
compared with the no-policy path under lower risk. Because lower risk is
itself a large change in the economy, I describe it first.

Using the same inherited households makes the initial comparison a response
to policy, rather than a difference between populations that have already
adjusted to it. The cohort measures then show whether the initial response
persists over the remaining fertile years. Finally, the path of household
entry records the consequences for the housing market after the additional
births reach adulthood. These measures answer different questions and need
not rank policies in the same order.

\subsection{Lower earnings risk}
\label{sec:policy-risk}

Figure~\ref{fig:policy-risk} compares the no-policy paths at the estimated and
at lower earnings risk. Lower risk raises births in 2023--2026 by 11.3 percent
and the completed fertility of the cohort aged 18 to 21 in 2023 by 16.0
percent, from 1.64 to 1.91. The effect is smaller than at a fixed price
(Section~\ref{sec:results-risk}) because the additional demand for family
space raises the house price, by 1.3 percent on impact and by about 10 percent
by 2063. For comparison, \citet{sommerFertilityChoiceLife2016} finds that moving from 1990s to 1970s
earnings risk raises completed fertility by 16 percent, and removing risk by
26 percent.

\problem{Lower earnings risk also compresses earnings inequality, so the
fertility response does not isolate the effect of uncertainty.}
I lower the standard deviation of the shocks by mixing the
transition matrix toward a less dispersed one with the same one-step expected
earnings in each state, and then scale all earnings states by a common factor
so that mean productivity stays at one in the new long-run distribution. On
impact the scaling lowers every household's earnings by 4.6 percent, which
works against births. Over time, the cross-sectional dispersion of earnings
falls, so lower risk here combines lower uncertainty with lower inequality,
and births rise steeply with earnings at the bottom of the distribution. Holding the earnings distribution at each age fixed while making part of
earnings heterogeneity known at entry would help distinguish uncertainty
from inequality.

\begin{figure}[t]
\centering
\includegraphics[width=0.95\linewidth]{\NT/risk0_baselines.pdf}
\caption{No-policy paths at the estimated and at lower earnings risk}
\label{fig:policy-risk}
\par\smallskip
\begin{minipage}{0.92\linewidth}
\footnotesize\textit{Notes:} Period fertility, completed fertility by cohort
(cohort aged 18--21 in the year shown), household heads and the house price
(2023 = 100), from 2023. Solid: estimated earnings risk ($\sigma_z=0.48$).
Dashed: lower earnings risk from 2023 on ($\sigma_z=0.35$), with mean
efficiency normalized in its long-run distribution.
\end{minipage}
\end{figure}

\subsection{Housing supply and mortgage credit}
\label{sec:policy-supply-credit}

Table~\ref{tab:policy-main} summarizes the results, and
Figures~\ref{fig:policy-impact} to~\ref{fig:policy-long} show them. Each
row in the table reports the effect relative to its own no-policy path.
The lower-risk rows therefore measure the additional effect of a housing
policy after reducing risk; they do not include the fertility increase from
lower risk itself. The contrast between impact births and completed fertility
is central: an early increase can disappear when later births are counted.

\begin{table}[t]
\centering
\caption{Policies from 2023: effects relative to the no-policy path}
\label{tab:policy-main}
\small
\begin{tabularx}{\textwidth}{@{}Xrrrrr@{}}
\toprule
 & Price & Births & \multicolumn{2}{c}{Completed fertility} & Households \\
 \cmidrule(lr){4-5}
 & 2023 & 2023--26 & cohort 2023 & cohort 2043 & 2063 \\
\midrule
\multicolumn{6}{@{}l}{\emph{Estimated earnings risk}}\\
Housing supply $+10\%$ & $-8.9$ & $+3.1$ & $+4.8$ & $+3.8$ & $+1.4$ \\
Mortgage credit, LTV 100 & $+0.5$ & $+1.9$ & $-0.4$ & $+0.7$ & $+0.2$ \\
Matched supply $+6.3\%$ & $-5.8$ & $+2.0$ & $+3.1$ & $+2.5$ & $+0.9$ \\
Credit for parents, LTV 100 & $+0.4$ & $+6.0$ & $+6.3$ & $+7.9$ & $+2.3$ \\
\addlinespace
\multicolumn{6}{@{}l}{\emph{Lower earnings risk}}\\
Housing supply $+10\%$ & $-9.1$ & $+2.6$ & $+4.0$ & $+3.3$ & $+1.3$ \\
Mortgage credit, LTV 100 & $+1.2$ & $+3.4$ & $-0.6$ & $-0.0$ & $+0.3$ \\
Matched supply $+13.1\%^{\dagger}$ & $-11.6$ & $+3.4$ & $+5.1$ & $+4.3$ & $+1.7$ \\
Credit for parents, LTV 100 & $+0.9$ & $+12.1$ & $+11.5$ & $+10.6$ & $+4.5$ \\
\bottomrule
\end{tabularx}
\par\smallskip
\begin{minipage}{0.95\linewidth}
\footnotesize\textit{Notes:} Percent differences from the no-policy path at
the same earnings risk, from the same 2023 state (estimated earnings risk: the
constant-decay transition). Completed fertility: children
ever born to the cohort aged 18--21 in the year shown. $^{\dagger}$ Under lower
risk the matched shift is sized to match credit's effect on births in
2023--2026 at that risk.
\end{minipage}
\end{table}
\problem{The lower-risk rows are from the earlier two-step transition; rerun
at the constant-decay state with the rescaled earnings.}

\paragraph{On impact.} Figure~\ref{fig:policy-impact} shows the effects on
the price and on births in 2023--2026. A 10 percent expansion of the housing
supply lowers the house price by 8.9 percent and raises births by 3.1 percent.
Full mortgage credit leaves the price nearly unchanged, $+0.5$ percent, and
raises births by 1.9 percent. The ranking depends on risk. Under lower risk,
supply raises births by 2.6 percent and credit by 3.4 percent: lower risk
shrinks the precautionary cost of a child that cheaper space relieves, as
\eqref{eq:r-risk-cost} predicts, and it raises the number of young households
held back only by the down payment. Births per percentage point of price
decline are about 0.35 on impact for supply.

The earnings profiles in the right-hand panels place these aggregate
responses alongside the household margin. Households in different earnings
deciles differ in the resources available for a child, so an aggregate
birth response averages over very different initial birth rates. Because
the figure groups households by current earnings, the profiles also include
differences in age, wealth and family composition within each decile. They
are not the response to changing the earnings of an otherwise identical
household.

\begin{figure}[t]
\centering
\includegraphics[width=\linewidth]{\NT/seq1_impact_risk_4panel_v2.pdf}
\caption{Housing supply and mortgage credit: on impact}
\label{fig:policy-impact}
\par\smallskip
\begin{minipage}{0.92\linewidth}
\footnotesize\textit{Notes:} Left: house price and births in 2023--2026,
percent difference from the no-policy path at the same earnings risk. Right:
births per household by decile of current labor earnings at fertile ages,
2023. Top: housing supply $+10\%$. Bottom: mortgage credit, LTV 80 to 100.
\end{minipage}
\end{figure}

\paragraph{Cohorts.} Figure~\ref{fig:policy-cohorts} shows completed
fertility by cohort. Here the two policies differ in kind. Supply raises the
completed fertility of the cohort aged 18 to 21 in 2023 by 4.8 percent, and of
every later cohort, by 3.8 percent for the 2043 cohort. Credit lowers the
completed fertility of the 2023 cohort by 0.4 percent and raises that of the
2043 cohort by 0.7 percent. The births credit adds on impact are mostly births
brought forward: for the 2023 cohort, credit raises birth rates at ages 18 to
21 and lowers them at every age after 26, and it lowers the mean age at first
birth by 0.12 years. Cohorts already past their early twenties in 2023 also
gain from both policies, since they still have births to have.

A cohort comparison holds the group of households fixed as it ages. It can
therefore distinguish a change in the number of children from a shift in
when those children are born. By contrast, the impact measure adds births
across all fertile ages in the first period of the policy. Pulling births
forward from later periods raises that measure even when the same
households ultimately have no additional children.

\begin{figure}[t]
\centering
\includegraphics[width=\linewidth]{\NT/cohorts_from_fertile_2023/seq2_cohorts_risk_2panel_from1999.pdf}
\caption{Housing supply and mortgage credit: completed fertility by cohort}
\label{fig:policy-cohorts}
\par\smallskip
\begin{minipage}{0.92\linewidth}
\footnotesize\textit{Notes:} Completed fertility by cohort, percent
difference from the no-policy path at the same earnings risk, from the oldest
cohort still fertile in 2023 (aged 42--45 then). The horizontal axis is the
year the cohort is aged 18--21; the dotted line marks the first cohort wholly
after 2023. Left: housing supply $+10\%$. Right: mortgage credit, LTV 80 to
100.
\end{minipage}
\end{figure}

\paragraph{The long view.} Figure~\ref{fig:policy-long} follows the two
policies to 2100. Supply keeps the price below its no-policy path for decades
and raises period fertility by about 0.06 births per woman by 2063. The
number of households rises slowly as the larger cohorts enter adulthood.
Credit raises period fertility little after the first decade. Household
numbers respond later than births because entry is tied to earlier birth
cohorts. Thus a small initial change in the number of households is
consistent with a persistent fertility response; it is not evidence that
the demographic channel is absent. As larger cohorts enter, their housing
demand also offsets part of the initial decline in the price under the
supply policy.

\begin{figure}[t]
\centering
\includegraphics[width=\linewidth]{\NT/seq3_transition_risk_full_cut2100.pdf}
\caption{Housing supply and mortgage credit: the long view}
\label{fig:policy-long}
\par\smallskip
\begin{minipage}{0.92\linewidth}
\footnotesize\textit{Notes:} House price, period fertility and household
heads, each policy against its own no-policy path at the same earnings risk,
2023--2100. Top: housing supply $+10\%$. Bottom: mortgage credit, LTV 80 to
100.
\end{minipage}
\end{figure}

%OLD-20261016 \subsection{Credit and the starter home}
\subsection{Credit, timing and tenure choice}
\label{sec:policy-starter}

Easier credit mainly lets young households buy the space they already rent.
Figure~\ref{fig:policy-tenure} shows households aged 22 to 29 at the start of
2023 after the choice. Without the policy, 49 percent of them own; with
full credit, 73 percent do, an increase of 24 percentage points, or 50
percent. Few of the new owners move into family-size homes, and most of the
increase in ownership is not accompanied by more space. Credit does not
change the relative value of a child, so it mostly moves the timing of births,
not their number.
\problem{The builder of Figure~\ref{fig:policy-tenure} is not in the
repository and the decomposition into rented space and new space is read off
the figure; add the construction and the share of new family-size owners who
are childless.}

\begin{figure}[t]
\centering
\includegraphics[width=\linewidth]{../../output/model/experiments/transitions_2023/decay_ship_20261014/figures/slide33_credit_ownership_space_decay.pdf}
\caption{Credit, ownership and space among young households}
\label{fig:policy-tenure}
\par\smallskip
\begin{minipage}{0.92\linewidth}
\footnotesize\textit{Notes:} Households aged 22--29 in 2023 after the choice,
estimated earnings risk, constant-decay state. Mortgage credit, LTV 80 to 100.
\end{minipage}
\end{figure}

Why does credit raise ownership without raising completed fertility?
Figure~\ref{fig:policy-starter}, computed on the earlier two-step path, shows
what households do with it, at the
start of 2027--2030. Without a down payment, young households buy early and
buy small. Among households aged 22 to 25, ownership rises from 46 to 68
percent, and the share owning a two-room home from 8 to 25 percent. Births per
100 households at those ages are 22.2 with and without the policy. A small
home accommodates a family's space requirement, but households that buy one early
are not the households whose fertility was held back.

The same pattern holds later in the transition. In 2059, mortgage credit
raises ownership among households aged 22 to 37 by 21 percentage points;
young owners live in smaller homes, 5.64 rooms against 6.36, and carry more
debt, 0.68 of the value of their home against 0.55. A supply expansion instead
gives young owners larger homes, 6.65 rooms, with the same leverage. Credit
buys earlier, smaller and more leveraged ownership; supply buys more space.

This difference explains why ownership alone is an incomplete measure of
housing support for families. A household can become an owner by replacing
a rented dwelling with a small purchased home while taking on more debt.
The tenure change then says little about the space it can afford over the
years children remain at home. The room and debt outcomes are necessary to
interpret the large ownership response.

The PSID offers a partial check on this mechanism. Among childless first-time
buyers, those who bought homes of four rooms or fewer were 14 percentage
points less likely to have a first birth within six years (standard error
6.7). The gap goes in the model's direction, but without an instrument it is
as consistent with selection, couples not planning children buying small, as
with an effect of the home. The model also places far more young owners in
small homes than the data do, so the size of this channel is likely
overstated.

\begin{figure}[t]
\centering
\includegraphics[width=\linewidth]{\NT/starter_home/starter_home_trap_by_age_7p9167_base.pdf}
\caption{Credit and the starter home}
\label{fig:policy-starter}
\par\smallskip
\begin{minipage}{0.92\linewidth}
\footnotesize\textit{Notes:} Households at the start of 2027--2030, by age,
estimated earnings risk. Left: homeowners, percent of households. Middle:
households owning a two-room home. Right: births per 100 households in
2027--2030. Grey: no policy. Orange: mortgage credit, LTV 100.
\end{minipage}
\end{figure}

\subsection{Timing versus level: matched supply}
\label{sec:policy-matched}

Demographers call \emph{quantum} how many children are born and \emph{tempo}
when \citep{bongaartsFeeneyQuantumTempo1998}. Figure~\ref{fig:policy-qt} sets
credit against a supply expansion in the 2023 economy, with the supply
expansion sized to raise births on impact by as much as credit does: 2.0 and
1.9 percent. The two paths then separate. Under credit the gain is nearly
gone by 2035, 0.2 percent, though a smaller lasting effect builds to 2.5
percent in the long run, with the number of households 2.5 percent higher and
the price 3.3 percent higher. Under supply, births are higher in every period,
by 5.2 percent in 2079 and by 6.3 percent in the long run, when the number of
households is 6.3 percent higher and the price is back at its no-policy level.
Credit is mostly tempo, supply is quantum.

\begin{figure}[t]
\centering
\includegraphics[width=\linewidth]{../../output/model/experiments/transitions_2023/decay_ship_20261014/figures/test_quantum_tempo/propB_births_households_two_panels.pdf}
\caption{Quantum and tempo: matched supply and mortgage credit}
\label{fig:policy-qt}
\par\smallskip
\begin{minipage}{0.92\linewidth}
\footnotesize\textit{Notes:} Births and household heads, percent difference
from the no-policy path, constant-decay state, estimated earnings risk. Credit:
LTV limit 80 to 100 (financed share $\phi=1$) in 2023. Supply: the supply curve
shifts out by 6.3 percent in 2023.
\end{minipage}
\end{figure}

The comparison so far mixes the size of the two policies with their nature. To
separate them, I shift the supply curve by 6.3 percent, an amount chosen so
that births in 2023--2026 rise by 2.0 percent, as under full credit (1.9
percent). Table~\ref{tab:policy-tempo} compares the 2023 cohort under the two
policies. With equal births on impact, matched supply raises completed
fertility by 3.1 percent and credit lowers it by 0.6 percent in the earlier
path (0.4 percent at the constant-decay state).
\problem{Recompute Table~\ref{tab:policy-tempo} at the constant-decay state.} Supply raises
the share of the cohort that ever has a child, by 2.4 percent, and raises
birth rates at every age. Credit lowers that share slightly and moves births
earlier. Supply is not neutral for timing either: it lowers the mean age at
first birth by 0.14 years, about as much as credit does for this cohort. A
cheaper room lowers the per-period cost of the commitment, so households
commit sooner as well as more often. Credit changes almost only the timing.

\begin{table}[t]
\centering
\caption{Matched supply and mortgage credit: the 2023 cohort (two-step path)}
\label{tab:policy-tempo}
\small
\begin{tabularx}{\textwidth}{@{}Xrrrr@{}}
\toprule
 & Completed & Ever had a & Mean age at & Mean age \\
 & fertility & first birth & first birth & at birth \\
\midrule
No policy & 1.641 & 0.676 & 27.10 & 31.99 \\
Mortgage credit, LTV 100 & 1.632 ($-0.6\%$) & 0.670 ($-0.8\%$) & 26.99 & 31.91 \\
Matched supply $+6.3\%$ & 1.692 ($+3.1\%$) & 0.692 ($+2.4\%$) & 26.96 & 31.92 \\
\bottomrule
\end{tabularx}
\par\smallskip
\begin{minipage}{0.95\linewidth}
\footnotesize\textit{Notes:} Cohort aged 18--21 in 2023, estimated earnings
risk. Rates are per surviving member of the cohort; mean ages are
rate-weighted midpoints of four-year age cells. These are descriptive cohort
sums, not a formal tempo adjustment \citep{bongaartsFeeneyQuantumTempo1998}.
\end{minipage}
\end{table}

The matched shift is a construct: its size is set by credit's effect on
impact, not by any proposed policy. Its purpose is to show that, for the same
number of births in the first four years, cheaper space raises the number of
children a cohort has and easier credit does not. It is not a comparison at
equal fiscal cost or equal welfare cost. The matching removes one dimension,
the initial birth response, so that the subsequent cohort paths can reveal
the difference between earlier births and additional births.

\subsection{Credit for parents}
\label{sec:policy-parents}

If credit fails because it reaches households whose fertility was not held
back, credit tied to children should do better. Credit available only once
there is a child lowers the cost of the child, not the cost of waiting. I raise
the financed share to 1 only for households with a child at home.
Figures~\ref{fig:policy-parents-ltrun} and~\ref{fig:policy-parents} show the
result. Births in 2023--2026 rise by 6.0 percent, three times the gain from
general credit (1.9 percent), and the completed fertility of the 2023 cohort by
6.3 percent, against $-0.4$ percent, with the price nearly unchanged on impact,
$+0.4$ percent. In the long run the number of households is 26 percent higher
with credit for parents and 2.5 percent higher with general credit. Under
lower earnings risk the effects roughly double, to 12.1 and 11.5 percent.

\begin{figure}[t]
\centering
\includegraphics[width=\linewidth]{../../output/model/experiments/transitions_2023/decay_ship_20261014/figures/test_quantum_tempo/credit_children_impact_longrun_decay.pdf}
\caption{Credit for parents: impact and long run}
\label{fig:policy-parents-ltrun}
\par\smallskip
\begin{minipage}{0.92\linewidth}
\footnotesize\textit{Notes:} Constant-decay state, estimated earnings risk,
percent difference from the no-policy path. Credit for parents: LTV 100 for
households with a child at home. General credit: LTV 100 for all households.
\end{minipage}
\end{figure}

The policy works through two channels that the experiment does not separate.
First, a birth unlocks ownership without a down payment, so it rewards the
birth itself. Second, a parent who owns can borrow up to the full value of the
home in every period with a child at home, which is a buffer in exactly the
state that carries the commitment. The second channel should weaken when risk
is lower, as earnings insurance does; the effect instead grows, which points
to the first: lower risk raises the demand to own, so access to ownership tied
to a child is worth more. A policy that relaxes the down payment only in the
period of the birth would separate the two.

\problem{The model does not price the default or mortgage-insurance cost
of credit for parents.}
The effect also depends on young households' demand for ownership, which the
model overstates (Section~\ref{sec:quant-2023}).

\begin{figure}[t]
\centering
\includegraphics[width=\linewidth]{\NT/pphi_block/pphi100_impact_and_cohorts.pdf}
\caption{Credit for parents: on impact and by cohort}
\label{fig:policy-parents}
\par\smallskip
\begin{minipage}{0.92\linewidth}
\footnotesize\textit{Notes:} LTV 100 for households with a child at home.
Left: house price and births in 2023--2026. Right: completed fertility by
cohort. Each against the no-policy path at the same earnings risk.
\end{minipage}
\end{figure}

Figure~\ref{fig:policy-parents-long} follows this response beyond the first
cohorts. The price effect grows as the larger birth cohorts enter the housing
market, and the number of households separates gradually from the no-policy
path. Period fertility remains higher over the horizon shown, although the
gap narrows later while housing becomes more expensive. The lower-risk case
has larger price and population responses as well as a larger initial birth
response. These paths are consistent with the housing-demand feedback in the
model, although they do not separately identify how much that feedback
limits the fertility response to credit for parents.

\begin{figure}[t]
\centering
\includegraphics[width=\linewidth]{\NT/pphi_block/pphi100_long_view_cut2100.pdf}
\caption{Credit for parents: the long view}
\label{fig:policy-parents-long}
\par\smallskip
\begin{minipage}{0.92\linewidth}
\footnotesize\textit{Notes:} House price, period fertility and household
heads against the no-policy path at the same earnings risk, 2023--2100.
\end{minipage}
\end{figure}

\subsection{The long run}
\label{sec:policy-long-run}

In the long run, births are pinned at replacement: the population is
stationary only if each household has 2.1 children on average. A policy that
raises desired fertility at a given price therefore shows up in the long run
as a higher price and more households, not as more births per household. The
house price is the one at which households choose replacement fertility, and
the number of households is the one that fills the housing stock at that
price. This is a statement about the stationary endpoint under the model's
renewal normalization. It does not imply that the transition has no gains in
completed fertility: cohorts can have more children for many periods while
the population adjusts toward its new level.

This logic gives a sharp prediction for supply. A 10 percent shift in the
supply curve leaves the replacement price unchanged, so in the long run the
price returns to its no-policy level and the number of households rises by 10
percent: the additional rooms are filled by additional families. The model
reproduces this exactly. For mortgage credit, the long-run price is 3.3
percent higher and the number of households 2.5 percent higher at the
estimated risk, so credit raises desired fertility slightly at a given price.
Under lower risk the long-run effect of credit is about zero and cannot be
signed at the precision of the solution. Credit for parents raises the
long-run price by 27 percent and the number of households by 26 percent.
Lower earnings risk alone raises the long-run price by 49 percent: births are
not very sensitive to the price of space in the model, so a large rise
in the price is needed to bring them back to replacement.

The long-run effect of credit also depends on the price level. Holding the
price at the low long-run level of the no-policy path, removing the down
payment raises births 1.0 percentage point above replacement; at prices about
50 percent higher, it raises them by only 0.1 to 0.2 percentage points. When space is expensive,
the per-period cost of a family room dominates the decision, and credit, which
changes only the one-time down payment, matters little.
% [MODEL] evaluation.md 2F: long-run base risk price +3.33%, households
% +2.51%, endpoint birth error 4e-7; lower risk endpoint accepted at a kink,
% price effect between -1.3% and -0.6%; fixed-price tests at P = 0.3533; endpoint
% scans over P in [0.515, 0.528]: +0.09 to +0.18 pt (mean +0.14). 2I: PPHI100
% long-run price 0.4496 vs 0.3533 (+27.2%), population 0.6378 vs 0.5072
% (+25.7%). Lower risk long run: housing_behavior table, price +49.2%.
